\documentclass[10pt,a4paper]{amsart}
\usepackage{amsmath,amssymb,amsthm,graphicx}
% Retained preamble for the Stacks Project chapter divisors.tex (Divisors).
% The chapter file pinned-original-source.tex carries only \input{preamble} on its first line, so no
% preamble text lives inside the chapter itself.  The Stacks Project uses one single shared preamble for
% every chapter; the snapshot stages it verbatim in sources/stacks-chow/preamble.tex.  The block below is
% copied mechanically, line for line, from that file's lines 14-18 (the Xy-pic setup) and lines 154-217
% (the native theorem declarations together with the mathematical macros).  Nothing was added, removed or
% redefined.  The wrapper class and the amsmath/amssymb/amsthm/graphicx base packages are supplied by the
% assembler, so the class/\externaldocument part of the shared preamble is not repeated here.
% For commutative diagrams we use Xy-pic
\usepackage[all]{xy}

% We use 2cell for 2-commutative diagrams.
\xyoption{2cell}
% Theorem environments.
%
\theoremstyle{plain}
\newtheorem{theorem}[subsection]{Theorem}
\newtheorem{proposition}[subsection]{Proposition}
\newtheorem{lemma}[subsection]{Lemma}

\theoremstyle{definition}
\newtheorem{definition}[subsection]{Definition}
\newtheorem{example}[subsection]{Example}
\newtheorem{exercise}[subsection]{Exercise}
\newtheorem{situation}[subsection]{Situation}

\theoremstyle{remark}
\newtheorem{remark}[subsection]{Remark}
\newtheorem{remarks}[subsection]{Remarks}

\numberwithin{equation}{subsection}

% Macros
%
\def\lim{\mathop{\mathrm{lim}}\nolimits}
\def\colim{\mathop{\mathrm{colim}}\nolimits}
\def\Spec{\mathop{\mathrm{Spec}}}
\def\Hom{\mathop{\mathrm{Hom}}\nolimits}
\def\Ext{\mathop{\mathrm{Ext}}\nolimits}
\def\SheafHom{\mathop{\mathcal{H}\!\mathit{om}}\nolimits}
\def\SheafExt{\mathop{\mathcal{E}\!\mathit{xt}}\nolimits}
\def\Sch{\mathit{Sch}}
\def\Mor{\mathop{\mathrm{Mor}}\nolimits}
\def\Ob{\mathop{\mathrm{Ob}}\nolimits}
\def\Sh{\mathop{\mathit{Sh}}\nolimits}
\def\NL{\mathop{N\!L}\nolimits}
\def\CH{\mathop{\mathrm{CH}}\nolimits}
\def\proetale{{pro\text{-}\acute{e}tale}}
\def\etale{{\acute{e}tale}}
\def\QCoh{\mathit{QCoh}}
\def\Ker{\mathop{\mathrm{Ker}}}
\def\Im{\mathop{\mathrm{Im}}}
\def\Coker{\mathop{\mathrm{Coker}}}
\def\Coim{\mathop{\mathrm{Coim}}}

% Boxtimes
%
\DeclareMathSymbol{\boxtimes}{\mathbin}{AMSa}{"02}

%
% Macros for moduli stacks/spaces
%
\def\QCohstack{\mathcal{QC}\!\mathit{oh}}
\def\Cohstack{\mathcal{C}\!\mathit{oh}}
\def\Spacesstack{\mathcal{S}\!\mathit{paces}}
\def\Quotfunctor{\mathrm{Quot}}
\def\Hilbfunctor{\mathrm{Hilb}}
\def\Curvesstack{\mathcal{C}\!\mathit{urves}}
\def\Polarizedstack{\mathcal{P}\!\mathit{olarized}}
\def\Complexesstack{\mathcal{C}\!\mathit{omplexes}}
% \Pic is the operator that assigns to X its picard group, usage \Pic(X)
% \Picardstack_{X/B} denotes the Picard stack of X over B
% \Picardfunctor_{X/B} denotes the Picard functor of X over B
\def\Pic{\mathop{\mathrm{Pic}}\nolimits}
\def\Picardstack{\mathcal{P}\!\mathit{ic}}
\def\Picardfunctor{\mathrm{Pic}}
\def\Deformationcategory{\mathcal{D}\!\mathit{ef}}

% --- appended 2026-10-12 (Chief): the project own definitions of the reference/slogan/history
% display environments, copied mechanically from the shared Stacks preamble lines 8-12
% (\usepackage{verbatim} plus the three \newenvironment lines, which the chapter preamble keeps
% verbatim from the project preamble). Nothing added, removed or redefined: these environments
% carry no mathematical content and the project itself typesets them as comments. Without these
% three lines a counted statement whose own text uses them fails with Environment slogan undefined.
\usepackage{verbatim}
\newenvironment{reference}{\comment}{\endcomment}
\newenvironment{slogan}{\comment}{\endcomment}
\newenvironment{history}{\comment}{\endcomment}
\usepackage{cleveref}
\crefname{originallemma}{original lemma}{original lemmas}
\Crefname{originallemma}{Original lemma}{Original lemmas}
\makeatletter
\newlabel{lemma-blow-up-effective-Cartier-divisor}{{0807}{source}{}{sourcepointer-1}{}}
\newlabel{lemma-blow-up-effective-Cartier-divisor@cref}{{[originallemma][0][]0807}{[1][0][]source}{}{}{}}
\makeatother
\crefname{originallemma}{original lemma}{original lemmas}
\Crefname{originallemma}{Original lemma}{Original lemmas}
\makeatletter
\newlabel{lemma-blowing-up-affine}{{0804}{source}{}{sourcepointer-2}{}}
\newlabel{lemma-blowing-up-affine@cref}{{[originallemma][0][]0804}{[1][0][]source}{}{}{}}
\makeatother
\crefname{originallemma}{original lemma}{original lemmas}
\Crefname{originallemma}{Original lemma}{Original lemmas}
\makeatletter
\newlabel{lemma-blowing-up-gives-effective-Cartier-divisor}{{02OS}{source}{}{sourcepointer-3}{}}
\newlabel{lemma-blowing-up-gives-effective-Cartier-divisor@cref}{{[originallemma][0][]02OS}{[1][0][]source}{}{}{}}
\makeatother
\crefname{originallemma}{original lemma}{original lemmas}
\Crefname{originallemma}{Original lemma}{Original lemmas}
\makeatletter
\newlabel{lemma-blowing-up-two-ideals}{{080A}{source}{}{sourcepointer-4}{}}
\newlabel{lemma-blowing-up-two-ideals@cref}{{[originallemma][0][]080A}{[1][0][]source}{}{}{}}
\makeatother
\crefname{originallemma}{original lemma}{original lemmas}
\Crefname{originallemma}{Original lemma}{Original lemmas}
\makeatletter
\newlabel{lemma-closed-subscheme-proj-finite-type}{{0803}{source}{}{sourcepointer-5}{}}
\newlabel{lemma-closed-subscheme-proj-finite-type@cref}{{[originallemma][0][]0803}{[1][0][]source}{}{}{}}
\makeatother
\crefname{originallemma}{original lemma}{original lemmas}
\Crefname{originallemma}{Original lemma}{Original lemmas}
\makeatletter
\newlabel{morphisms-lemma-closed-immersion-finite-presentation}{{01TV}{source}{}{sourcepointer-6}{}}
\newlabel{morphisms-lemma-closed-immersion-finite-presentation@cref}{{[originallemma][0][]01TV}{[1][0][]source}{}{}{}}
\makeatother
\crefname{originallemma}{original lemma}{original lemmas}
\Crefname{originallemma}{Original lemma}{Original lemmas}
\makeatletter
\newlabel{properties-lemma-quasi-coherent-finite-type-ideals}{{01PH}{source}{}{sourcepointer-7}{}}
\newlabel{properties-lemma-quasi-coherent-finite-type-ideals@cref}{{[originallemma][0][]01PH}{[1][0][]source}{}{}{}}
\makeatother
\title{The composition of two blowing ups in centres of finite presentation is again a blowing up}
\author{The Stacks Project Authors}
\date{stacks-project:divisors.tex:local-snapshot-20261012; GNU FDL 1.2 or later}
\begin{document}
\maketitle
\noindent\textit{Editorial note.} All delivered mathematical text is the author's own, copied from the pinned Stacks Project chapter file pinned-original-source.tex; no statement or proof step was rewritten, reordered or removed.\par\medskip
\noindent\textit{Editorial note.} Following the multi-claim convention for this source, source\_url is the source-level URL https://stacks.math.columbia.edu/ and the source label lemma-composition-finite-type-blowups is carried in claim\_key and here rather than in a per-claim tag URL.\par\medskip
\noindent\textit{Editorial note.} The preamble is the chapter's shared Stacks Project preamble, retained as sources/stacks-divisors/retained-preamble.tex: a line-for-line excerpt of the shared Stacks preamble staged with this snapshot. No macro was invented or redefined, and this card adds no preamble text of its own.\par\medskip
\noindent\textit{Editorial note.} Statement span 8656-8673 and proof span 8675-8740 come from the halfA pool card operations/pool-cards/stacks-divisors-halfA.json and were re-read in the pinned chapter file before this card was written.\par\medskip
\noindent\textit{Editorial note.} The statement carries an author slogan ('Composition of blowing ups is a blowing up.'). The project's own shared preamble, retained in full at 88 lines, defines the slogan environment as a comment environment, so the slogan is not typeset, exactly as in the project's own output; it is nevertheless part of the delivered statement span and is preserved verbatim.\par\medskip
\noindent\textit{Editorial note.} The proof uses the author's notation for the ideal of the exceptional divisor and for the ideal of the second centre, and it introduces the letter J for the submodule regarded as an ideal; all of this is preserved unchanged.\par\medskip
\noindent\textit{Editorial note.} 7 cross-references occur inside the counted spans; they are listed in external\_references as navigation pointers only.\par\medskip
\noindent\textit{Allowed standard prerequisites (editorial).}\begin{itemize}
\item A closed subscheme of finite presentation is cut out by a finite type quasi-coherent sheaf of ideals, and the complement of the support of such an ideal is a quasi-compact open, cited from the morphisms and properties chapters. Cited at proof lines 8675-8677 and 8680-8682. Supplies the ideal whose Rees algebra presents the first blowing up and makes the relevant complements quasi-compact.
\item The blowing up along a finite type ideal is an isomorphism away from the centre of the ideal, cited from the same chapter. Cited at proof lines 8682-8684. Identifies the preimage of the complement of the first centre and keeps the complement of the union of the two images quasi-compact.
\item The description of a closed subscheme of the blowing up as the blowing up of the base in a finite type submodule of a power of the defining ideal, supported outside a prescribed quasi-compact open, cited from the same chapter. Cited at proof lines 8687-8690. Is the step that replaces the second centre, which a priori lives on the blowing up, by a single ideal on the base.
\item The blowing up in a product of ideals equals the blowing up in the pulled back ideal, and blowing up in an invertible ideal, equivalently in an effective Cartier divisor, does not change the scheme, both cited from the same chapter. Cited at proof lines 8710-8712 and 8719-8722. Reduces the claim to the identification of two ideal sheaves on the second blowing up.
\item The local description of an affine blowing up and of its blowup algebra, cited from the same chapter. Cited at proof lines 8726-8728. Turns the ideal equality into the elementary computation with a quotient of the form m over a to the d.
\end{itemize}\medskip
\ifdefined\hypertarget\hypertarget{sourcepointer-1}{}\fi
\noindent\textit{Source reference.} Original lemma 0807: focused-seeds/AGP/sources/stacks-snapshot-20261011/tags.txt, source lines 10172--10172. Same-chapter reference used in the retained author text at source line(s) 8726, 8955. It is supplied only as a navigation pointer so that the retained text compiles; the counted item delivers the author's own statement and the author's complete proof, and no part of the referenced item is copied in or substituted for the counted argument.\par\medskip
\ifdefined\hypertarget\hypertarget{sourcepointer-2}{}\fi
\noindent\textit{Source reference.} Original lemma 0804: focused-seeds/AGP/sources/stacks-snapshot-20261011/tags.txt, source lines 10169--10169. Same-chapter reference used in the retained author text at source line(s) 8406, 8524, 8536, 8578, 8731, 8837, 8956, 9203, 9275. It is supplied only as a navigation pointer so that the retained text compiles; the counted item delivers the author's own statement and the author's complete proof, and no part of the referenced item is copied in or substituted for the counted argument.\par\medskip
\ifdefined\hypertarget\hypertarget{sourcepointer-3}{}\fi
\noindent\textit{Source reference.} Original lemma 02OS: focused-seeds/AGP/sources/stacks-snapshot-20261011/tags.txt, source lines 3527--3527. Same-chapter reference used in the retained author text at source line(s) 8444, 8552, 8600, 8685, 8765, 9106. It is supplied only as a navigation pointer so that the retained text compiles; the counted item delivers the author's own statement and the author's complete proof, and no part of the referenced item is copied in or substituted for the counted argument.\par\medskip
\ifdefined\hypertarget\hypertarget{sourcepointer-4}{}\fi
\noindent\textit{Source reference.} Original lemma 080A: focused-seeds/AGP/sources/stacks-snapshot-20261011/tags.txt, source lines 10175--10175. Same-chapter reference used in the retained author text at source line(s) 8717, 8727, 8954, 9167. It is supplied only as a navigation pointer so that the retained text compiles; the counted item delivers the author's own statement and the author's complete proof, and no part of the referenced item is copied in or substituted for the counted argument.\par\medskip
\ifdefined\hypertarget\hypertarget{sourcepointer-5}{}\fi
\noindent\textit{Source reference.} Original lemma 0803: focused-seeds/AGP/sources/stacks-snapshot-20261011/tags.txt, source lines 10168--10168. Same-chapter reference used in the retained author text at source line(s) 8689. It is supplied only as a navigation pointer so that the retained text compiles; the counted item delivers the author's own statement and the author's complete proof, and no part of the referenced item is copied in or substituted for the counted argument.\par\medskip
\ifdefined\hypertarget\hypertarget{sourcepointer-6}{}\fi
\noindent\textit{Source reference.} Original lemma 01TV: focused-seeds/AGP/sources/stacks-snapshot-20261011/tags.txt, source lines 2388--2388. Same-chapter reference used in the retained author text at source line(s) 1451, 5010, 5119, 8124, 8679, 9099. It is supplied only as a navigation pointer so that the retained text compiles; the counted item delivers the author's own statement and the author's complete proof, and no part of the referenced item is copied in or substituted for the counted argument.\par\medskip
\ifdefined\hypertarget\hypertarget{sourcepointer-7}{}\fi
\noindent\textit{Source reference.} Original lemma 01PH: focused-seeds/AGP/sources/stacks-snapshot-20261011/tags.txt, source lines 2230--2230. Same-chapter reference used in the retained author text at source line(s) 1453, 8683, 9187. It is supplied only as a navigation pointer so that the retained text compiles; the counted item delivers the author's own statement and the author's complete proof, and no part of the referenced item is copied in or substituted for the counted argument.\par\medskip
\begin{lemma}
\label{lemma-composition-finite-type-blowups}
\begin{slogan}
Composition of blowing ups is a blowing up
\end{slogan}
Let $X$ be a quasi-compact and quasi-separated scheme.
Let $Z \subset X$ be a closed subscheme of finite presentation.
Let $b : X' \to X$ be the blowing up with center $Z$. Let $Z' \subset X'$ be
a closed subscheme of finite presentation.
Let $X'' \to X'$ be the blowing up with center $Z'$.
There exists a closed subscheme $Y \subset X$ of finite presentation,
such that
\begin{enumerate}
\item $Y = Z \cup b(Z')$ set theoretically, and
\item the composition $X'' \to X$ is isomorphic to the blowing up
of $X$ in $Y$.
\end{enumerate}
\end{lemma}

\begin{proof}
The condition that $Z \to X$ is of finite presentation means that
$Z$ is cut out by a finite type quasi-coherent sheaf of ideals
$\mathcal{I} \subset \mathcal{O}_X$, see
Morphisms, Lemma \ref{morphisms-lemma-closed-immersion-finite-presentation}.
Write $\mathcal{A} = \bigoplus_{n \geq 0} \mathcal{I}^n$ so that
$X' = \underline{\text{Proj}}(\mathcal{A})$.
Note that $X \setminus Z$ is a quasi-compact open of $X$ by
Properties, Lemma \ref{properties-lemma-quasi-coherent-finite-type-ideals}.
Since $b^{-1}(X \setminus Z) \to X \setminus Z$ is an isomorphism
(Lemma \ref{lemma-blowing-up-gives-effective-Cartier-divisor}) the same
result shows that
$b^{-1}(X \setminus Z) \setminus Z'$ is quasi-compact open in $X'$.
Hence $U = X \setminus (Z \cup b(Z'))$ is quasi-compact open in $X$.
By Lemma \ref{lemma-closed-subscheme-proj-finite-type}
there exist a $d > 0$ and a finite type
$\mathcal{O}_X$-submodule $\mathcal{F} \subset \mathcal{I}^d$ such
that $Z' = \underline{\text{Proj}}(\mathcal{A}/\mathcal{F}\mathcal{A})$
and such that the support of $\mathcal{I}^d/\mathcal{F}$ is contained
in $X \setminus U$.

\medskip\noindent
Since $\mathcal{F} \subset \mathcal{I}^d$ is an $\mathcal{O}_X$-submodule
we may think of $\mathcal{F} \subset \mathcal{I}^d \subset \mathcal{O}_X$
as a finite type quasi-coherent sheaf of ideals on $X$. Let's denote this
$\mathcal{J} \subset \mathcal{O}_X$ to prevent confusion. Since
$\mathcal{I}^d / \mathcal{J}$ and $\mathcal{O}/\mathcal{I}^d$
are supported on $X \setminus U$ we see that $V(\mathcal{J})$ is contained
in $X \setminus U$. Conversely, as $\mathcal{J} \subset \mathcal{I}^d$
we see that $Z \subset V(\mathcal{J})$. Over
$X \setminus Z \cong X' \setminus b^{-1}(Z)$ the sheaf of ideals
$\mathcal{J}$ cuts out $Z'$ (see displayed formula below). Hence
$V(\mathcal{J})$ equals $Z \cup b(Z')$. It follows that also
$V(\mathcal{I}\mathcal{J}) = Z \cup b(Z')$ set theoretically. Moreover,
$\mathcal{I}\mathcal{J}$ is an ideal of finite type as a product of two such.
We claim that $X'' \to X$ is isomorphic to the blowing up of $X$ in
$\mathcal{I}\mathcal{J}$ which finishes the proof of the lemma by setting
$Y = V(\mathcal{I}\mathcal{J})$.

\medskip\noindent
First, recall that the blowup of $X$ in $\mathcal{I}\mathcal{J}$
is the same as the blowup of $X'$ in $b^{-1}\mathcal{J} \mathcal{O}_{X'}$,
see Lemma \ref{lemma-blowing-up-two-ideals}.
Hence it suffices to show that the blowup of $X'$ in
$b^{-1}\mathcal{J} \mathcal{O}_{X'}$ agrees with the blowup of $X'$
in $Z'$. We will show that
$$
b^{-1}\mathcal{J} \mathcal{O}_{X'} = \mathcal{I}_E^d \mathcal{I}_{Z'}
$$
as ideal sheaves on $X''$. This will prove what we want as
$\mathcal{I}_E^d$ cuts out the effective Cartier divisor $dE$
and we can use Lemmas \ref{lemma-blow-up-effective-Cartier-divisor} and
\ref{lemma-blowing-up-two-ideals}.

\medskip\noindent
To see the displayed equality of the ideals we may work locally.
With notation $A$, $I$, $a \in I$ as in Lemma \ref{lemma-blowing-up-affine}
we see that $\mathcal{F}$ corresponds to an $R$-submodule $M \subset I^d$
mapping isomorphically to an ideal $J \subset R$. The condition
$Z' = \underline{\text{Proj}}(\mathcal{A}/\mathcal{F}\mathcal{A})$
means that $Z' \cap \Spec(A[\frac{I}{a}])$ is cut out by the ideal
generated by the elements $m/a^d$, $m \in M$. Say the element $m \in M$
corresponds to the function $f \in J$. Then in the affine blowup algebra
$A' = A[\frac{I}{a}]$ we see that $f = (a^dm)/a^d = a^d (m/a^d)$.
Thus the equality holds.
\end{proof}

\end{document}
